\section{PAPER 026: Discovery Calibration — EMA-Based Community Rate Tracking and Contra\-diction Resolution for Autonomous KG Research}

\textbf{Author}: Bryan Alexander Buchorn  
\textbf{Affiliation}: Independent Researcher, Las Vegas, NV, USA  
\textbf{Status}: v2.52.0 (Phase 172 (STRB) COMPLETE)
\textbf{Date}: May 2, 2026

\textbf{CEREBRUM Phase 73}

---

\subsection{Abstract}

We present three Phase 73 components that address systematic biases in CEREBRUM's autonomous knowledge graph discovery pipeline. \textbf{DiscoveryCalibrator} tracks per-community scan and discovery rates via Exponential Moving Average (EMA) and applies an inverse-rate sampling multiplier to steer \texttt{ResearchAgent} toward under\-represented graph regions. \textbf{Contra\-dictionResolver} classifies candidate findings into four evidence classes (clean / revision\_candidate / contested / discardable) using a deter\-ministic Noisy-OR evidence model, filtering the most conflicted proposals before they reach \texttt{AutoApprover}. \textbf{CandidateRegistry} replaces the flat evaluated-pairs set with a TTL-aware registry that applies nomination-count boosts and enforces memory bounds via LRU eviction. Together, these components reduce redundant computation, prevent community saturation, and improve the precision of the discovery pipeline by pre-filtering struc\-turally contr\-adicted proposals.

---

\subsection{1. Motivation: Systematic Bias in Autonomous Discovery}

An uncal\-ibrated \texttt{ResearchAgent} exhibits two systematic failure modes:

\begin{enumerate}
\item \textbf{Community saturation}: The agent repeatedly scans the same densely-connected communities (high community weight) because they yield frequent high-confidence candidates. Underr\-epresented sparse communities receive no coverage, leaving genuine missing links undis\-covered.

\item \textbf{Contra\-diction blindness}: Proposals with strong contr\-adicting evidence in the existing graph (high \texttt{contradiction\_score}) are forwarded to \texttt{AutoApprover} and consume classifier compute, even when deter\-ministic evidence-weight analysis would immediately classify them as discardable.

\end{enumerate}
---

\subsection{2. DiscoveryCalibrator}

\subsubsection{2.1 Per-Community EMA Tracking}

For each community \texttt{c}, the calibrator maintains:
\begin{itemize}
\item \texttt{scan\_rate(c)}: EMA of scans per unit time
\item \texttt{discovery\_rate(c)}: EMA of approved discoveries per scan

\begin{lstlisting}
calibrator.record_scan(community_id)
calibrator.record_discovery(community_id)
\end{lstlisting}

\end{itemize}
EMA update: \texttt{rate\_t = α × event + (1 - α) × rate\_{t-1}} where \texttt{α = 0.1} (default).

\subsubsection{2.2 Inverse-Rate Sampling Multiplier}

The community weight \texttt{w(c)} used in \texttt{\_score\_discovery\_potential()}:

\begin{lstlisting}
w(c) = global_discovery_rate / (discovery_rate(c) + ε)
\end{lstlisting}

Cold-start communities (never scanned) receive \texttt{max\_weight = 5.0}. Communities with higher-than-global discovery rates receive \texttt{w \textless  1.0} (suppressed). Communities with lower rates receive \texttt{w \textgreater  1.0} (boosted).

\subsubsection{2.3 Temporal Recency Scoring}

\texttt{ValidationReport.recency\_score} is computed via exponential decay from the publication year:

\begin{lstlisting}
recency_score = exp(-λ × max(0, current_year - pub_year))
\end{lstlisting}

Default \texttt{λ} corresponds to a 7-year half-life. Recent literature (\textless  2 years old) scores $\ge$ 0.9.

---

\subsection{3. Contra\-dictionResolver}

\subsubsection{3.1 Evidence Model}

For a given finding with \texttt{proposed\_confidences = [c\_1, ..., c\_k]} and \texttt{contradiction\_score}:

\begin{lstlisting}
net_evidence_score = Noisy-OR(c_1, ..., c_k) - contradiction_score
\end{lstlisting}

Noisy-OR: \texttt{1 - ∏(1 - c\_i)} — the probability of at least one path being correct.

\subsubsection{3.2 Resolution Classes}

\begin{center}
{\footnotesize
\begin{tabularx}{\columnwidth}{|>{\raggedright\arraybackslash}X|>{\raggedright\arraybackslash}X|>{\raggedright\arraybackslash}X|}
\hline
net\_evidence\_score & Class & Action \\ \hline
$\ge$ 0.6 & \texttt{"clean"} & Forward to AutoApprover \\ \hline
0.3 – 0.6 & \texttt{"revision\_candidate"} & Queue in \texttt{\_revision\_candidates} \\ \hline
0.1 – 0.3 & \texttt{"contested"} & Forward to AutoApprover with penalty \\ \hline
\textless  0.1 & \texttt{"discardable"} & Auto-reject; never reaches AutoApprover \\ \hline
\end{tabularx}
}
\end{center}

Revision candidates accumulate in \texttt{ResearchAgent.\_revision\_candidates} for periodic batch review.

---

\subsection{4. CandidateRegistry}

\subsubsection{4.1 TTL-Aware Registry}

Replaces the flat \texttt{\_evaluated\_pairs: Set[Tuple[str, str]]} with a dict-based registry:

\begin{lstlisting}
registry[pair] = CandidateRecord(
    nomination_count=N,
    first_seen=t_0,
    last_seen=t_N,
    ttl=timedelta(hours=24),
)
\end{lstlisting}

\texttt{is\_registered(pair)} returns True and blocks re-evaluation if within TTL.

\subsubsection{4.2 Nomination Boost}

When a pair is nominated multiple times (by different scan cycles or reasoning strategies), its \texttt{discovery\_potential} receives a log-scale boost:

\begin{lstlisting}
boosted_potential = raw_potential × min(log(N + 1) + 1, max_boost)
\end{lstlisting}

Default \texttt{max\_boost = 3.0}. This rewards repeatedly-surfaced candidates without linearly amplifying noise.

\subsubsection{4.3 Memory Bound}

\texttt{max\_entries} (default 10,000) triggers LRU eviction: the registry evicts the least-recently-seen entry when capacity is reached, ensuring bounded memory usage in long-running deployments.

---

\subsection{5. References}

\begin{itemize}
\item [pearl\-2000\-causality] Pearl, J. (2000). \textit{Causality: Models, reasoning, and inference.} Cambridge University Press.

\end{itemize}
---
\textbf{Copyright © 2026 Bryan Alexander Buchorn. All Rights Reserved.}

---
\subsection{Acknow\-ledgments}

The author gratefully ackno\-wledges the use of Claude (Anthropic) as a research assistant throughout this work. Claude assisted with mathe\-matical forma\-lization, code generation, manuscript preparation, and technical writing. All conceptual contr\-ibutions, archi\-tectural decisions, exper\-imental design, and intel\-lectual claims are solely the author's.

\textbf{Reviewed on}: May 2, 2026 for version v2.52.0
